
\documentclass[10pt,landscape]{article}
\usepackage{mathwizards-formulario}

% ── Comandos propios de este documento ─────────
\renewcommand{\headrulewidth}{0pt}
\pagestyle{fancy}
\fancyhf{}
\fancyfoot[C]{\color{primary}\small\textbf{F\'{i}sica I -- Leyes del Movimiento | UDO N\'{u}cleo Monagas}}

\begin{document}
\pagecolor{white}

% Título principal
\begin{center}
{\Huge\bfseries\color{primary} FORMULARIO DE F\'{I}SICA I}\\[2pt]
{\large\color{darkbg}\textbf{Leyes del Movimiento -- Newton} | Hoja Carta Horizontal}
\end{center}
\vspace{-4pt}

\begin{multicols}{3}

% ======================== COLUMNA 1 ========================

\cajatitulo{1. LEYES DE NEWTON}

\subsection*{Primera Ley (Inercia)}
Un cuerpo permanece en reposo o en MRU a menos que una fuerza neta act\'{u}e sobre \'{e}l.
\[
\sum \vec{F} = 0 \;\Rightarrow\; \vec{a} = 0
\]

\subsection*{Segunda Ley}
La aceleraci\'{o}n es proporcional a la fuerza neta e inversamente proporcional a la masa.
\[
\boxed{\vec{F}_{\text{neta}} = m\vec{a}} \qquad \boxed{\sum F_x = ma_x} \qquad \boxed{\sum F_y = ma_y}
\]

\subsection*{Tercera Ley (Acci\'{o}n--Reacci\'{o}n)}
\[
\vec{F}_{12} = -\vec{F}_{21}
\]
Las fuerzas de acci\'{o}n y reacci\'{o}n son iguales en magnitud, opuestas en direcci\'{o}n, y act\'{u}an sobre cuerpos \textbf{diferentes}.

\vspace{4pt}
\cajatitulo{2. FUERZAS FUNDAMENTALES}

\subsection*{Peso (Fuerza Gravitacional)}
\[
\boxed{W = mg} \qquad g \approx 9.8\,\text{m/s}^2 \;(\text{cerca de la Tierra})
\]

\subsection*{Fuerza Normal ($N$)}
Perpendicular a la superficie de contacto. No siempre es igual al peso.
\begin{itemize}[leftmargin=*,nosep,topsep=0pt]
  \item Plano horizontal sin fuerzas verticales: $N = mg$
  \item Plano inclinado: $N = mg\cos\theta$
  \item Fuerza $F$ con \'{a}ngulo $\theta$ sobre horizontal: $N = mg - F\sin\theta$
\end{itemize}

\subsection*{Fuerza de Rozamiento}
\textbf{Est\'{a}tico:}
\[
\boxed{f_s \leq \mu_s N} \qquad f_{s,\text{m\'{a}x}} = \mu_s N
\]
\textbf{Cin\'{e}tico:}
\[
\boxed{f_k = \mu_k N} \qquad \text{(constante, opuesta al movimiento)}
\]
\cajainfo{\textbf{Nota:} $\mu_s > \mu_k$ generalmente. El rozamiento est\'{a}tico var\'{i}a entre $0$ y $\mu_s N$.}

\subsection*{Fuerza El\'{a}stica (Ley de Hooke)}
\[
\boxed{F = -kx} \qquad \boxed{T = 2\pi\sqrt{\frac{m}{k}}} \qquad \boxed{\omega = \sqrt{\frac{k}{m}}}
\]
Energ\'{i}a potencial el\'{a}stica: $U = \frac{1}{2}kx^2$

\vspace{4pt}
\cajatitulo{3. DIN\'{A}MICA CIRCULAR}

\subsection*{Aceleraci\'{o}n Centr\'{i}peta}
\[
\boxed{a_c = \frac{v^2}{r} = \omega^2 r} \qquad \boxed{\omega = \frac{2\pi}{T} = 2\pi f}
\]

\subsection*{Fuerza Centr\'{i}peta}
\[
\boxed{F_c = \frac{mv^2}{r} = m\omega^2 r}
\]

\subsection*{Peralte de Curvas (sin rozamiento lateral)}
\[
\boxed{\tan\theta = \frac{v^2}{rg}} \qquad v_{\text{m\'{a}x}} = \sqrt{rg\tan\theta}
\]

% ======================== COLUMNA 2 ========================

\cajatitulo{4. CANTIDAD DE MOVIMIENTO}

\subsection*{Definici\'{o}n}
\[
\boxed{\vec{p} = m\vec{v}} \qquad [\text{p}] = \text{kg}\cdot\text{m/s}
\]

\subsection*{Segunda Ley en t\'{e}rminos de $p$}
\[
\boxed{\vec{F} = \frac{d\vec{p}}{dt}}
\]

\subsection*{Principio de Conservaci\'{o}n}
Si $\sum \vec{F}_{\text{ext}} = 0$, entonces:
\[
\boxed{\vec{p}_{\text{total}} = \text{constante}} \qquad \sum m_i \vec{v}_i = \text{constante}
\]

\subsection*{Impulso}
\[
\boxed{\vec{J} = \vec{F}\Delta t = \Delta \vec{p} = m\vec{v}_f - m\vec{v}_i}
\]
Fuerza promedio: $\displaystyle F_{\text{prom}} = \frac{\Delta p}{\Delta t}$

\vspace{4pt}
\cajatitulo{5. COLISIONES}

\subsection*{Colisi\'{o}n El\'{a}stica (se conserva $p$ y $E_k$)}
\[
\sum \vec{p}_{\text{antes}} = \sum \vec{p}_{\text{despu\'{e}s}} \qquad \sum E_k^{\text{antes}} = \sum E_k^{\text{despu\'{e}s}}
\]
Para masas $m_1$, $m_2$ en 1D:
\[
v_{1f} = \frac{(m_1-m_2)v_{1i} + 2m_2 v_{2i}}{m_1+m_2}
\]

\subsection*{Colisi\'{o}n Inel\'{a}stica (se conserva solo $p$)}
\[
\sum \vec{p}_{\text{antes}} = \sum \vec{p}_{\text{despu\'{e}s}}
\]
\textbf{Perfectamente inel\'{a}stica} (se pegan):
\[
\boxed{m_1\vec{v}_1 + m_2\vec{v}_2 = (m_1+m_2)\vec{v}_f}
\]

\vspace{4pt}
\cajatitulo{6. PLANOS INCLINADOS}

\subsection*{Fuerzas sobre un bloque}
\begin{itemize}[leftmargin=*,nosep,topsep=0pt]
  \item Componente del peso paralela: $mg\sin\theta$
  \item Componente del peso perpendicular: $mg\cos\theta$
  \item Normal: $N = mg\cos\theta$
\end{itemize}

\subsection*{Aceleraci\'{o}n (con rozamiento cin\'{e}tico)}
\textbf{Bajando:}
\[
\boxed{a = g(\sin\theta - \mu_k \cos\theta)}
\]
\textbf{Subiendo:}
\[
\boxed{a = -g(\sin\theta + \mu_k \cos\theta)}
\]

\subsection*{\'{A}ngulo de reposo}
Cuando el bloque est\'{a} a punto de deslizarse:
\[
\boxed{\tan\theta_{\text{m\'{a}x}} = \mu_s}
\]

\vspace{4pt}
\cajatitulo{7. POLEAS Y SISTEMAS}

\subsection*{Dos masas conectadas (Atwood)}
\[
a = \frac{m_2 - m_1}{m_1 + m_2}g \qquad T = \frac{2m_1 m_2}{m_1 + m_2}g
\]

\subsection*{Masa sobre mesa + masa colgante}
Si $m_1$ est\'{a} sobre mesa sin fricci\'{o}n y $m_2$ cuelga:
\[
a = \frac{m_2}{m_1 + m_2}g \qquad T = m_1 a = \frac{m_1 m_2}{m_1 + m_2}g
\]

\subsection*{Sistema con fricci\'{o}n}
Incluir $f_k = \mu_k N$ en la direcci\'{o}n opuesta al movimiento en la ecuaci\'{o}n de Newton.

% ======================== COLUMNA 3 ========================

\cajatitulo{8. ENERG\'{I}A Y TRABAJO}

\subsection*{Trabajo}
\[
\boxed{W = \vec{F} \cdot \vec{d} = Fd\cos\theta} \qquad [W] = \text{J} = \text{N}\cdot\text{m}
\]

\subsection*{Energ\'{i}a Cin\'{e}tica}
\[
\boxed{K = \frac{1}{2}mv^2}
\]

\subsection*{Energ\'{i}a Potencial Gravitacional}
\[
\boxed{U_g = mgh}
\]

\subsection*{Teorema Trabajo--Energ\'{i}a}
\[
\boxed{W_{\text{neto}} = \Delta K = K_f - K_i}
\]

\subsection*{Conservaci\'{o}n de Energ\'{i}a Mec\'{a}nica}
\[
\boxed{K_i + U_i = K_f + U_f} \quad \text{(sin fuerzas no conservativas)}
\]

\vspace{4pt}
\cajatitulo{9. MOVIMIENTO ARM\'{O}NICO SIMPLE (MAS)}

\subsection*{Ecuaciones del MAS}
\[
x(t) = A\cos(\omega t + \phi) \qquad v(t) = -A\omega\sin(\omega t + \phi)
\]
\[
a(t) = -A\omega^2\cos(\omega t + \phi) = -\omega^2 x(t)
\]

\subsection*{Frecuencia y Periodo}
\[
\boxed{\omega = \sqrt{\frac{k}{m}}} \qquad \boxed{T = 2\pi\sqrt{\frac{m}{k}}} \qquad \boxed{f = \frac{1}{2\pi}\sqrt{\frac{k}{m}}}
\]

\subsection*{Fuerza en el MAS}
\[
\boxed{F = -kx = ma = -m\omega^2 x}
\]
M\'{a}xima fuerza: $F_{\text{m\'{a}x}} = kA = m\omega^2 A$ (en los extremos)

\vspace{4pt}
\cajatitulo{10. CA\'{I}DA LIBRE Y PROYECTILES}

\subsection*{Ca\'{i}da Libre (sin resistencia)}
\[
\boxed{v_f^2 = v_i^2 + 2gh} \qquad \boxed{v_f = v_i + gt} \qquad \boxed{h = v_i t + \frac{1}{2}gt^2}
\]
Fuerza neta: $F_{\text{neta}} = mg$ (hacia abajo)

\subsection*{Con resistencia del aire $R$ constante}
\[
F_{\text{neta}} = mg - R \quad \Rightarrow \quad a = g - \frac{R}{m}
\]

\subsection*{Movimiento Parab\'{o}lico}
\[
x = v_0\cos\theta \cdot t \qquad y = v_0\sin\theta \cdot t - \frac{1}{2}gt^2
\]
Alcance m\'{a}ximo: $\displaystyle R = \frac{v_0^2\sin(2\theta)}{g}$

\vspace{4pt}
\cajatitulo{11. CONSEJOS DE RESOLUCI\'{O}N}

\begin{enumerate}[leftmargin=*,nosep,topsep=0pt]
  \item \textbf{DCL siempre:} Dibujar el diagrama de cuerpo libre para cada cuerpo.
  \item \textbf{Ejes inteligentes:} Un eje paralelo a la aceleraci\'{o}n y otro perpendicular.
  \item \textbf{Signos consistentes:} Definir positivo en la direcci\'{o}n del movimiento esperado.
  \item \textbf{Fricci\'{o}n:} Verificar si hay movimiento ($f_k$) o est\'{a} en reposo ($f_s \leq \mu_s N$).
  \item \textbf{Cuerdas ideales:} Tensi\'{o}n constante a lo largo de la cuerda, masa despreciable.
  \item \textbf{Poleas ideales:} Cambian direcci\'{o}n de la tensi\'{o}n, sin fricci\'{o}n ni masa.
  \item \textbf{Conservaci\'{o}n:} Verificar si hay fuerzas externas netas antes de aplicar conservaci\'{o}n de $p$.
  \item \textbf{Unidades:} Trabajar en SI: kg, m, s, N. Convertir g a kg, cm a m.
\end{enumerate}

\vspace{4pt}
\cajatitulo{12. CONSTANTES \'{U}TILES}

\begin{center}
\begin{tabular}{@{}ll@{}}
\toprule
\textbf{Constante} & \textbf{Valor} \\
\midrule
$g$ (gravedad) & $9.8\,\text{m/s}^2$ \\
$G$ (gravitacional) & $6.67 \times 10^{-11}\,\text{N}\cdot\text{m}^2/\text{kg}^2$ \\
Masa Tierra & $5.97 \times 10^{24}\,\text{kg}$ \\
Radio Tierra & $6.37 \times 10^6\,\text{m}$ \\
\bottomrule
\end{tabular}
\end{center}

\end{multicols}

\end{document}
